\documentclass[12pt]{article}
\usepackage[margin=1in]{geometry}
\usepackage[utf8]{inputenc}
\usepackage[english]{babel}
\usepackage{setspace}
\usepackage{url}
\usepackage{footmisc}
\usepackage{fancyhdr}

\doublespacing

\pagestyle{fancy}
\fancyhf{}
\fancyhead[R]{Vaught \thepage}  % Top right corner: "Vaught" and page number

\begin{document}


\bigskip

The \textit{Hippocratic Oath}, often attributed to Hippocrates\footnote{Reference 1 states that there is evidence that the Oath was written after Hippocrates death},
has been recited for nearly 2,500 years by physicians as they enter the medical profession\footnote{Reference 2 states that reciting the oath became more common approx. 500 years ago}. 
Despite its venerable history, many contemporary ethicists have dismissed this ancient pledge as outmoded\footnote{Reference 3 points out issues within the Classical Oath today}, 
citing massive transformations in science, economics, and politics since the era of its creation. 
Yet, a closer reading of the oath suggests that it still speaks forcefully to modern medicine’s core dilemmas: 
the balance between intervention and restraint, the distinction between healing and harm, and the moral imperative to serve patients and society.

\bigskip

The oath begins with an appeal to “Apollo the physician and [by] Asclepius and Hygieia and Panacea.”\footnote{Reference 4, Classic Hippocratic Oath}
Although many can and will decry this invocation of Greek gods, one might reflect on the deeper symbolism involved. Apollo embodied the healing arts
\footnote{Reference 5 contains a breakdown of Apollo's place in Greek culture}, while Asclepius’s name signified ``to give well-being"
\footnote{Reference 6 shows that his name may not have had this meaning, but simply a mistranslation}, reminding physicians of their sacred duty to 
heal. In a modern context, this religious dimension could be replaced by a personal or universal moral commitment, yet the underlying message remains: 
medicine demands humility before higher ethical standards and a devotion to patient welfare. A second contested segment of the oath is its prohibition 
on providing “a drug that is deadly to anyone if asked,”\footnote{Reference 4, Classic Hippocratic Oath} commonly interpreted as barring euthanasia. 
This line likely addressed the ancient Greek fear of physicians enlisted to assassinate political enemies.\footnote{Reference 7 states that physicians 
assisting in assassinations was common, thus this conclusion is logical} Scholars note that Hippocrates likely had no specific intention to weigh in on 
the concept of assisted suicide or “euthanasia,” a term not coined until the nineteenth century\footnote{Reference 8 shows that the first usage of 
`euthanasia' was in 1869}. Nevertheless, the oath’s refusal to condone deliberate killing can be read as a stern warning against abusing medical 
authority - whether for political, punitive, or malicious ends. Such caution surely endures today in discussions about torture, executions, and 
any situation where a doctor might be coerced into causing harm contrary to ethical practice.

\bigskip

Likewise, the oath’s perceived ban on surgery - “I will cede this [surgery] to men who are practitioners of this activity”
\footnote{Reference 4, Classic Hippocratic Oath} - has prompted further criticism. However, historical data confirm that 
ancient Greek practitioners did perform invasive procedures, sometimes quite aggressively.\footnote{Reference 9 shows a history of surgery, and shows Greek history of advanced surgery} The text, more accurately, 
seems to acknowledge professional specialization: a generalist who lacks adequate training should defer to those with proven expertise. 
This foreshadowed modern multidisciplinary teamwork, in which radiotherapists, oncologists, and surgeons each address specific 
patient needs. Far from hindering innovation, the oath advises humility, collaboration, and recognition of one’s limits. 
Objections have also arisen around the line “I will not give a woman a destructive pessary,”
\footnote{Reference 4, Classic Hippocratic Oath} interpreted by some as a prohibition of abortion. Closer examination reveals 
that ancient Greece allowed abortions in various forms, particularly oral methods.\footnote{Reference 10 shows abortion through history, includes a section on Greek abortion} The caution likely targeted the lethal 
infections associated with certain pessary-based interventions rather than a blanket ban on abortion. Thus, the oath’s 
stance can be read as a plea for safe medical procedures, urging physicians to avoid reckless or harmful treatments - an ethos 
that remains pertinent in the modern debate over ensuring safe reproductive care.

\bigskip

Beyond these procedural concerns, the most crucial elements of the oath may be its strident admonitions regarding harm and injustice.
From counseling practitioners to avoid sexual or other forms of exploitation to its strong advocacy of physician responsibility for the vulnerable, 
the oath anticipates long-standing issues in medical ethics. In the present age, where the cost of health care and the presence of financial conflicts of interest\footnote{Reference 11 shows pharmacies increasing prices, resulting in 7.3Bn in profit } loom large, the vow to put patient needs first resonates
just as powerfully. Indeed, the oath’s directive to combat personal and social injustice underscores a physician’s duty to advocate for universal
health care, affordable medicines, and policies that mitigate inequities in access. High costs and restricted coverage have generated suffering 
for millions of Americans, evidenced by delays in treatment, medical bankruptcies, and a general sense of frustration among both patients and 
providers. While the Affordable Care Act improved coverage for certain groups, structural problems persist: the system still excludes many, 
and physicians grapple with administrative burdens, diminishing autonomy, and constant cost pressures. In this context, the 
Hippocratic Oath’s message against harm and injustice can be read as urging doctors to take a more active role in championing reform - 
whether through direct patient advocacy or broader policy engagement.

\bigskip

Although penned in an era radically different from our own, the \textit{Hippocratic Oath} should continue to challenge 
physicians to navigate the complexities of modern health care with moral clarity. 
Its references to Greek gods reflect an abiding call to humility and compassion, its stance on euthanasia and surgery underscores the 
judicious exercise of medical authority, and its injunction against injustice remains a prompt to tackle systemic inequities. 
Rather than viewing the oath as an outdated relic, one might recognize its enduring relevance in addressing debates over physician autonomy, 
universal care, and the ethical obligations that come with medical expertise. 
If anything, the Oath’s spirit is needed more than ever: “I swear by what is most sacred, my patients.”\footnote{Reference 4, Classic Hippocratic Oath}

\bigskip

\newpage
\begin{thebibliography}{9}

\bibitem{historyguild}
History Guild. "The Hippocratic Oath: Changing Perspectives." Accessed January 20, 2025.  
\url{https://historyguild.org/the-hippocratic-oath-changing-perspectives/}

\bibitem{einsteinmed}
Hulkower, Raphael. "The History of the Hippocratic Oath: Outdated, Inauthentic, and Yet Still Relevant." \textit{Albert Einstein College of Medicine}, Bronx, New York 10461. Accessed January 20, 2025.  
\url{https://einsteinmed.edu/uploadedFiles/EJBM/page41_page44.pdf}

\bibitem{springer}
SpringerLink. "Why the Hippocratic Oath Is Still Relevant." Accessed January 20, 2025.  
\url{https://link.springer.com/article/10.1007/s11017-011-9203-z}

\bibitem{hippocratic_oath}
Edelstein, Ludwig, translator. \textit{Hippocratic Oath.} Baltimore: Johns Hopkins University Press, 1943. Accessed via PBS website, January 20, 2025.  
\url{https://www.pbs.org/wgbh/nova/article/hippocratic-oath-today/}

\bibitem{apollo}
Theoi Project. "Apollon." Accessed January 20, 2025.  
\url{https://www.theoi.com/Olympios/Apollon.html}

\bibitem{asclepius}
Mythopedia. "Asclepius." Accessed January 20, 2025.  
\url{https://mythopedia.com/topics/asclepius}

\bibitem{oath_relevance}
Baker Institute. "Deciphering the Hippocratic Oath in the 21st Century." Accessed January 20, 2025.  
\url{https://www.bakerinstitute.org/research/why-hippocratic-oath-still-relevant}

\bibitem{euthanasia}
Oxford English Dictionary. "Euthanasia, n." Accessed January 20, 2025.  
\url{https://www.oed.com/dictionary/euthanasia_n}

\bibitem{surgery}
Wasim Shah Mushtque Shah et al. "History of Surgery in Ancient, Middle and Modern Era – A Review Article." \textit{International Journal of Creative Research Thoughts (IJCRT)}, Volume 9, Issue 11, November 2021. Accessed January 20, 2025.  
\url{https://ijcrt.org/papers/IJCRT2111305.pdf}

\bibitem{abortion_ancient}
Garroway, Kristine Henriksen. "Abortion in the Ancient Near East and Greco-Roman World." \textit{TheTorah.com}, 2023. Accessed January 20, 2025.  
\url{https://www.thetorah.com/article/abortion-in-the-ancient-near-east-and-greco-roman-world}

\bibitem{finance}
Financial Times. "US Pharmacy Groups Marked Up Drug Prices for \$7.3bn Gain, Regulator Says." Accessed January 20, 2025.  
\url{https://www.ft.com/content/318bef8a-88fe-495a-ba4c-917fb133be6b}

\end{thebibliography}

\end{document}